\documentclass[12pt,a4paper]{article}

% --- Codificación e Idioma ---
\usepackage[utf8]{inputenc}
\usepackage[T1]{fontenc}
\usepackage[spanish,es-noshorthands]{babel}

% --- Geometría ---
\usepackage[a4paper, margin=2.5cm]{geometry}

% --- Matemáticas y Unidades ---
\usepackage{amsmath}
\usepackage[separate-uncertainty=true, locale=DE]{siunitx}

% --- Tablas ---
\usepackage{booktabs}
\usepackage{array}

% --- Listas ---
\usepackage{enumitem}

% --- Diagramas TikZ ---
\usepackage{tikz}
\usetikzlibrary{positioning, arrows}

% --- Hipervínculos ---
\usepackage{hyperref}
\hypersetup{
    colorlinks=true,
    linkcolor=blue,
    urlcolor=cyan,
    citecolor=blue,
}

% --- Metadatos ---
\title{Análisis Comparativo de Chips Medidores de Energía\\de Analog Devices}
\author{}
\date{16 de junio de 2026}

\begin{document}

\maketitle

\section{Introducción}

Los circuitos integrados (ICs) \textbf{ADE7753}, \textbf{ADE7758} y \textbf{ADE7880},
fabricados por \textbf{Analog Devices (ADI)}, representan un estándar en la industria
para la medición de energía eléctrica. Se utilizan ampliamente en el diseño de
medidores inteligentes (\textit{smart meters}), analizadores de calidad de la energía y
sistemas de gestión energética industrial.

A continuación, se presenta un análisis detallado de cada uno de ellos.

\section{ADE7753 --- Medición Monofásica}

Chip de alto rendimiento diseñado para aplicaciones de medición de energía en sistemas
monofásicos.

\begin{itemize}[itemsep=1ex]
    \item \textbf{Fase:} Monofásico (\textit{single-phase}).
    \item \textbf{Funciones clave:} Mide energía activa, reactiva y aparente. También
          calcula los valores RMS de tensión y corriente.
    \item \textbf{Característica distintiva:} Integrador digital interno
          ($\frac{di}{dt}$) seleccionable, que permite interactuar directamente con
          sensores de corriente basados en bobinas de Rogowski, además de
          transformadores de corriente (CT) o resistencias \textit{shunt}.
    \item \textbf{Comunicación:} Interfaz serial SPI.
\end{itemize}

\section{ADE7758 --- Medición Trifásica Estándar}

El ADE7758 es la versión polifásica de la familia, escalando las capacidades de medición
a sistemas trifásicos.

\begin{itemize}[itemsep=1ex]
    \item \textbf{Fase:} Polifásico / Trifásico (\textit{polyphase}). Compatible con
          sistemas de 3 o 4 hilos (estrella o delta).
    \item \textbf{Funciones clave:} Mismo conjunto de mediciones que el ADE7753 (energía
          activa, reactiva, aparente, V-RMS, I-RMS) para cada una de las tres fases.
    \item \textbf{Característica distintiva:} Ampliamente utilizado en medidores
          trifásicos comerciales e industriales. Dispone de salidas de pulsos para
          calibración y monitoreo de energía.
    \item \textbf{Comunicación:} Interfaz serial SPI.
\end{itemize}

\section{ADE7880 --- Medición Trifásica Avanzada}

Chip moderno y potente enfocado en aplicaciones de alta gama que requieren no solo medir
el consumo, sino también analizar la calidad de la energía.

\begin{itemize}[itemsep=1ex]
    \item \textbf{Fase:} Polifásico / Trifásico.
    \item \textbf{Funciones clave:} Mediciones de energía y RMS con alta precisión.
    \item \textbf{Característica distintiva:} \textbf{Análisis de armónicos} mediante un
          motor de cálculo dedicado que extrae información de armónicos individuales de
          tensión y corriente (hasta el armónico 63). Ideal para analizadores de calidad
          de energía.
    \item \textbf{Comunicación:} SPI, I2C e interfaz de captura de datos de alta
          velocidad (HSDC).
\end{itemize}

\newpage
\section{Diagrama de Bloques Típico}

La Figura~\ref{fig:diagrama} muestra la arquitectura típica de conexión de estos
medidores. Las señales analógicas de tensión y corriente se acondicionan antes de
ingresar a los pines ADC del chip. Mediante un bus SPI o I2C, un microcontrolador
externo lee los registros internos del ADE.

\begin{figure}[h!]
    \centering
    \begin{tikzpicture}[
        node distance=1.5cm and 2cm,
        block/.style={rectangle, draw, minimum width=2.4cm, minimum height=1cm,
                      align=center, fill=blue!10, rounded corners},
        sensor/.style={rectangle, draw, minimum width=2cm, minimum height=0.8cm,
                       align=center, fill=green!10},
        comm/.style={rectangle, draw, minimum width=2cm, minimum height=0.8cm,
                     align=center, fill=orange!10},
        mcu/.style={rectangle, draw, minimum width=2.4cm, minimum height=1cm,
                    align=center, fill=red!10, rounded corners},
        arrow/.style={-stealth, thick},
    ]

        \node[sensor] (vsen) {Sensor\\de Tensión};
        \node[sensor, below=of vsen] (csen) {Sensor\\de Corriente};

        \node[block, right=2.5cm of vsen, minimum height=3.2cm] (ade)
              {ADE7753/58/80\\{\footnotesize ADC + DSP}};

        \node[mcu, right=of ade] (mcu) {Micro-\\controlador};

        \node[comm, below=of mcu] (comm) {SPI / I2C};

        \draw[arrow] (vsen.east) -- ++(1,0) |- ([yshift=4mm] ade.west);
        \draw[arrow] (csen.east) -- ++(1,0) |- ([yshift=-4mm] ade.west);
        \draw[arrow] (ade.east) -- (mcu.west);
        \draw[arrow] (mcu.south) -- (comm.north);

        \node[above, font=\small] at ([yshift=2mm] vsen.east) {V};
        \node[above, font=\small] at ([yshift=2mm] csen.east) {I};

        \node[comm, above=of ade] (ps) {Fuente\\de Alimentación};
        \draw[arrow] (ps.south) -- (ade.north);

    \end{tikzpicture}
    \caption{Diagrama de bloques típico de conexión de un chip ADE con un
             microcontrolador.}
    \label{fig:diagrama}
\end{figure}

\section{Tabla Comparativa}

\begin{table}[h!]
    \centering
    \caption{Comparación de características clave.}
    \renewcommand{\arraystretch}{1.6}
    \begin{tabular}{@{}l >{\raggedright\arraybackslash}p{3.2cm}
                        >{\raggedright\arraybackslash}p{3.5cm}
                        >{\raggedright\arraybackslash}p{3.5cm}@{}}
        \toprule
        \textbf{Característica} & \textbf{ADE7753} & \textbf{ADE7758}
                                & \textbf{ADE7880} \\
        \midrule
        \textbf{Fase}     & Monofásico & Trifásico & Trifásico \\
        \addlinespace
        \textbf{Funciones base}
            & Energía Act./React./Aparente,\newline V-RMS, I-RMS
            & Energía Act./React./Aparente,\newline V-RMS, I-RMS (por fase)
            & Energía Act./React./Aparente,\newline V-RMS, I-RMS (por fase) \\
        \addlinespace
        \textbf{Característica única}
            & Integrador digital $\frac{di}{dt}$ para bobinas de Rogowski
            & Estándar industrial para medidores trifásicos
            & \textbf{Análisis de armónicos} (hasta el 63) \\
        \addlinespace
        \textbf{Comunicación} & SPI & SPI & SPI, I2C, HSDC \\
        \addlinespace
        \textbf{Resolución ADC} & 20 bits & 20 bits & 24 bits \\
        \addlinespace
        \textbf{Alimentación} & \SI{5}{\V} & \SI{5}{\V} & \SI{3.3}{\V} \\
        \bottomrule
    \end{tabular}
    \label{tab:comparativa}
\end{table}

\section{Conclusión}

La elección entre estos tres circuitos integrados depende directamente de la aplicación:

\begin{itemize}
    \item \textbf{ADE7753:} Ideal para medidores monofásicos residenciales o
          sub-medidores, especialmente si se planea usar sensores de corriente
          flexibles como las bobinas de Rogowski.
    \item \textbf{ADE7758:} Opción robusta y probada para medidores trifásicos
          estándar en entornos comerciales o industriales donde no se requiere un
          análisis profundo de la calidad de la energía.
    \item \textbf{ADE7880:} Elección predilecta para equipos de gama alta como
          analizadores de calidad de la red, sistemas de protección de subestaciones
          o en investigación, donde el análisis de armónicos es un requisito
          fundamental.
\end{itemize}

\begin{thebibliography}{3}
    \bibitem{ade7753}
    Analog Devices, \textit{ADE7753 --- Single-Phase Active and Apparent Energy
    Metering IC}, Datasheet Rev.~E, 2011.
    \href{https://www.analog.com/media/en/technical-documentation/data-sheets/ADE7753.pdf}
         {Enlace al datasheet}.

    \bibitem{ade7758}
    Analog Devices, \textit{ADE7758 --- Polyphase Multifunction Energy Metering IC},
    Datasheet Rev.~C, 2010.
    \href{https://www.analog.com/media/en/technical-documentation/data-sheets/ADE7758.pdf}
         {Enlace al datasheet}.

    \bibitem{ade7880}
    Analog Devices, \textit{ADE7880 --- Polyphase Multifunction Energy Metering IC
    with Harmonic Analysis}, Datasheet Rev.~A, 2012.
    \href{https://www.analog.com/media/en/technical-documentation/data-sheets/ADE7880.pdf}
         {Enlace al datasheet}.
\end{thebibliography}

\end{document}
